\documentclass[11pt,a4paper]{article}

% shared/macros.tex
% Shared notation for all surface-partition math documents.
% Include with: % shared/macros.tex
% Shared notation for all surface-partition math documents.
% Include with: % shared/macros.tex
% Shared notation for all surface-partition math documents.
% Include with: \input{../shared/macros}

% -------------------------------------------------------------------------
% Packages (include here so each document only needs \input{../shared/macros})
% -------------------------------------------------------------------------
\usepackage{amsmath,amssymb,amsthm}
\usepackage{bm}           % \bm for bold math
\usepackage{mathtools}    % \coloneqq, dcases, etc.
\usepackage{booktabs}     % \toprule etc. in tables
\usepackage{hyperref}
\usepackage{geometry}
\usepackage{microtype}
\usepackage{parskip}
\usepackage{xcolor}
\usepackage{tcolorbox}
\usepackage{enumitem}
\usepackage{float}             % [H] float placement

\geometry{a4paper, margin=2.5cm}

% -------------------------------------------------------------------------
% Theorem-like environments
% -------------------------------------------------------------------------
\theoremstyle{definition}
\newtheorem{defn}{Definition}[section]
\newtheorem{remark}{Remark}[section]

\theoremstyle{plain}
\newtheorem{prop}{Proposition}[section]

% -------------------------------------------------------------------------
% Mesh and geometry
% -------------------------------------------------------------------------
\newcommand{\V}{\mathbf{V}}            % vertex array V[i] in R^dim
\newcommand{\F}{\mathbf{F}}            % face (triangle) array
\newcommand{\vone}[1]{v_{1,#1}}       % first vertex index of VP i
\newcommand{\vtwo}[1]{v_{2,#1}}       % second vertex index of VP i

% -------------------------------------------------------------------------
% Variable points
% -------------------------------------------------------------------------
\newcommand{\lam}{\lambda}             % scalar VP parameter
\newcommand{\Lam}{\boldsymbol{\lambda}}% full parameter vector
\newcommand{\vp}[1]{\mathbf{p}_{#1}}  % 3D position of VP i
\newcommand{\ed}[1]{\mathbf{d}_{#1}}  % edge direction d_i = V[v1_i] - V[v2_i]

% -------------------------------------------------------------------------
% Steiner / triple-point
% -------------------------------------------------------------------------
\newcommand{\Sv}{\mathbf{S}}           % Steiner (Fermat-Torricelli) point
\newcommand{\Ki}{K_{i}}               % projection matrix (I-n_i n_i^T)/r_i
\newcommand{\Mmat}{M}                 % sum of K_i matrices at a triple point

% -------------------------------------------------------------------------
% Normals and unit vectors
% -------------------------------------------------------------------------
\newcommand{\nhat}{\hat{\mathbf{n}}}  % unit normal to a triangle
\newcommand{\Dhat}{\hat{\boldsymbol{\Delta}}}  % unit segment direction

% -------------------------------------------------------------------------
% Operations
% -------------------------------------------------------------------------
\newcommand{\norm}[1]{\left\|#1\right\|}
\newcommand{\abs}[1]{\left|#1\right|}
\newcommand{\ip}[2]{\langle #1,\, #2 \rangle}

% Projection perpendicular to a unit vector \uhat
\newcommand{\Proj}[1]{\bigl(\mathbf{I} - #1 #1^\top\bigr)}

% argmax operator (winner-take-all assignment in Phase 1 documents)
\DeclareMathOperator*{\argmax}{arg\,max}
\DeclareMathOperator*{\argmin}{arg\,min}

% -------------------------------------------------------------------------
% Calculus shorthands
% -------------------------------------------------------------------------
\newcommand{\pd}[2]{\dfrac{\partial #1}{\partial #2}}
\newcommand{\pdd}[3]{\dfrac{\partial^2 #1}{\partial #2\,\partial #3}}
\newcommand{\grad}{\nabla}
\newcommand{\Hess}[1]{\nabla^2 #1}

% -------------------------------------------------------------------------
% Optimization problem objects
% -------------------------------------------------------------------------
\newcommand{\Preg}{P_{\mathrm{reg}}}  % regular perimeter
\newcommand{\Pst}{P_{\mathrm{st}}}   % Steiner perimeter
\newcommand{\Areg}{A^{\mathrm{reg}}} % regular cell area
\newcommand{\Ast}{A^{\mathrm{st}}}   % Steiner area contribution
\newcommand{\Atgt}{\bar{A}}           % target area per cell
\newcommand{\Lagr}{\mathcal{L}}       % Lagrangian
\newcommand{\objfac}{\sigma}          % IPOPT objective scaling factor

% -------------------------------------------------------------------------
% Code cross-references (for "In the code..." callouts)
% -------------------------------------------------------------------------
\newcommand{\code}[1]{\texttt{#1}}
\newcommand{\codefile}[1]{\texttt{\small #1}}

% -------------------------------------------------------------------------
% Threshold dynamics / approach B (10-mbo-auction-dynamics)
% -------------------------------------------------------------------------
\newcommand{\Dlump}{D}                 % lumped mass matrix diag(v)
\newcommand{\Atau}{A_\tau}             % discrete diffusion operator (M + tau K)^{-1} M
\newcommand{\Etau}{E_\tau}             % heat-content (threshold-dynamics) energy
\newcommand{\Rcell}{R_{\mathrm{cell}}} % radius of the equal-area disc sqrt(Abar/pi)
\newcommand{\hmean}{h_{\mathrm{mean}}} % mean edge length
\newcommand{\hmax}{h_{\mathrm{max}}}   % max edge length
\newcommand{\Gt}{G_t}                  % heat kernel at time t
\newcommand{\Ktau}{K^{\mathrm{res}}_\tau} % resolvent (backward-Euler) kernel


% -------------------------------------------------------------------------
% Packages (include here so each document only needs % shared/macros.tex
% Shared notation for all surface-partition math documents.
% Include with: \input{../shared/macros}

% -------------------------------------------------------------------------
% Packages (include here so each document only needs \input{../shared/macros})
% -------------------------------------------------------------------------
\usepackage{amsmath,amssymb,amsthm}
\usepackage{bm}           % \bm for bold math
\usepackage{mathtools}    % \coloneqq, dcases, etc.
\usepackage{booktabs}     % \toprule etc. in tables
\usepackage{hyperref}
\usepackage{geometry}
\usepackage{microtype}
\usepackage{parskip}
\usepackage{xcolor}
\usepackage{tcolorbox}
\usepackage{enumitem}
\usepackage{float}             % [H] float placement

\geometry{a4paper, margin=2.5cm}

% -------------------------------------------------------------------------
% Theorem-like environments
% -------------------------------------------------------------------------
\theoremstyle{definition}
\newtheorem{defn}{Definition}[section]
\newtheorem{remark}{Remark}[section]

\theoremstyle{plain}
\newtheorem{prop}{Proposition}[section]

% -------------------------------------------------------------------------
% Mesh and geometry
% -------------------------------------------------------------------------
\newcommand{\V}{\mathbf{V}}            % vertex array V[i] in R^dim
\newcommand{\F}{\mathbf{F}}            % face (triangle) array
\newcommand{\vone}[1]{v_{1,#1}}       % first vertex index of VP i
\newcommand{\vtwo}[1]{v_{2,#1}}       % second vertex index of VP i

% -------------------------------------------------------------------------
% Variable points
% -------------------------------------------------------------------------
\newcommand{\lam}{\lambda}             % scalar VP parameter
\newcommand{\Lam}{\boldsymbol{\lambda}}% full parameter vector
\newcommand{\vp}[1]{\mathbf{p}_{#1}}  % 3D position of VP i
\newcommand{\ed}[1]{\mathbf{d}_{#1}}  % edge direction d_i = V[v1_i] - V[v2_i]

% -------------------------------------------------------------------------
% Steiner / triple-point
% -------------------------------------------------------------------------
\newcommand{\Sv}{\mathbf{S}}           % Steiner (Fermat-Torricelli) point
\newcommand{\Ki}{K_{i}}               % projection matrix (I-n_i n_i^T)/r_i
\newcommand{\Mmat}{M}                 % sum of K_i matrices at a triple point

% -------------------------------------------------------------------------
% Normals and unit vectors
% -------------------------------------------------------------------------
\newcommand{\nhat}{\hat{\mathbf{n}}}  % unit normal to a triangle
\newcommand{\Dhat}{\hat{\boldsymbol{\Delta}}}  % unit segment direction

% -------------------------------------------------------------------------
% Operations
% -------------------------------------------------------------------------
\newcommand{\norm}[1]{\left\|#1\right\|}
\newcommand{\abs}[1]{\left|#1\right|}
\newcommand{\ip}[2]{\langle #1,\, #2 \rangle}

% Projection perpendicular to a unit vector \uhat
\newcommand{\Proj}[1]{\bigl(\mathbf{I} - #1 #1^\top\bigr)}

% argmax operator (winner-take-all assignment in Phase 1 documents)
\DeclareMathOperator*{\argmax}{arg\,max}
\DeclareMathOperator*{\argmin}{arg\,min}

% -------------------------------------------------------------------------
% Calculus shorthands
% -------------------------------------------------------------------------
\newcommand{\pd}[2]{\dfrac{\partial #1}{\partial #2}}
\newcommand{\pdd}[3]{\dfrac{\partial^2 #1}{\partial #2\,\partial #3}}
\newcommand{\grad}{\nabla}
\newcommand{\Hess}[1]{\nabla^2 #1}

% -------------------------------------------------------------------------
% Optimization problem objects
% -------------------------------------------------------------------------
\newcommand{\Preg}{P_{\mathrm{reg}}}  % regular perimeter
\newcommand{\Pst}{P_{\mathrm{st}}}   % Steiner perimeter
\newcommand{\Areg}{A^{\mathrm{reg}}} % regular cell area
\newcommand{\Ast}{A^{\mathrm{st}}}   % Steiner area contribution
\newcommand{\Atgt}{\bar{A}}           % target area per cell
\newcommand{\Lagr}{\mathcal{L}}       % Lagrangian
\newcommand{\objfac}{\sigma}          % IPOPT objective scaling factor

% -------------------------------------------------------------------------
% Code cross-references (for "In the code..." callouts)
% -------------------------------------------------------------------------
\newcommand{\code}[1]{\texttt{#1}}
\newcommand{\codefile}[1]{\texttt{\small #1}}

% -------------------------------------------------------------------------
% Threshold dynamics / approach B (10-mbo-auction-dynamics)
% -------------------------------------------------------------------------
\newcommand{\Dlump}{D}                 % lumped mass matrix diag(v)
\newcommand{\Atau}{A_\tau}             % discrete diffusion operator (M + tau K)^{-1} M
\newcommand{\Etau}{E_\tau}             % heat-content (threshold-dynamics) energy
\newcommand{\Rcell}{R_{\mathrm{cell}}} % radius of the equal-area disc sqrt(Abar/pi)
\newcommand{\hmean}{h_{\mathrm{mean}}} % mean edge length
\newcommand{\hmax}{h_{\mathrm{max}}}   % max edge length
\newcommand{\Gt}{G_t}                  % heat kernel at time t
\newcommand{\Ktau}{K^{\mathrm{res}}_\tau} % resolvent (backward-Euler) kernel
)
% -------------------------------------------------------------------------
\usepackage{amsmath,amssymb,amsthm}
\usepackage{bm}           % \bm for bold math
\usepackage{mathtools}    % \coloneqq, dcases, etc.
\usepackage{booktabs}     % \toprule etc. in tables
\usepackage{hyperref}
\usepackage{geometry}
\usepackage{microtype}
\usepackage{parskip}
\usepackage{xcolor}
\usepackage{tcolorbox}
\usepackage{enumitem}
\usepackage{float}             % [H] float placement

\geometry{a4paper, margin=2.5cm}

% -------------------------------------------------------------------------
% Theorem-like environments
% -------------------------------------------------------------------------
\theoremstyle{definition}
\newtheorem{defn}{Definition}[section]
\newtheorem{remark}{Remark}[section]

\theoremstyle{plain}
\newtheorem{prop}{Proposition}[section]

% -------------------------------------------------------------------------
% Mesh and geometry
% -------------------------------------------------------------------------
\newcommand{\V}{\mathbf{V}}            % vertex array V[i] in R^dim
\newcommand{\F}{\mathbf{F}}            % face (triangle) array
\newcommand{\vone}[1]{v_{1,#1}}       % first vertex index of VP i
\newcommand{\vtwo}[1]{v_{2,#1}}       % second vertex index of VP i

% -------------------------------------------------------------------------
% Variable points
% -------------------------------------------------------------------------
\newcommand{\lam}{\lambda}             % scalar VP parameter
\newcommand{\Lam}{\boldsymbol{\lambda}}% full parameter vector
\newcommand{\vp}[1]{\mathbf{p}_{#1}}  % 3D position of VP i
\newcommand{\ed}[1]{\mathbf{d}_{#1}}  % edge direction d_i = V[v1_i] - V[v2_i]

% -------------------------------------------------------------------------
% Steiner / triple-point
% -------------------------------------------------------------------------
\newcommand{\Sv}{\mathbf{S}}           % Steiner (Fermat-Torricelli) point
\newcommand{\Ki}{K_{i}}               % projection matrix (I-n_i n_i^T)/r_i
\newcommand{\Mmat}{M}                 % sum of K_i matrices at a triple point

% -------------------------------------------------------------------------
% Normals and unit vectors
% -------------------------------------------------------------------------
\newcommand{\nhat}{\hat{\mathbf{n}}}  % unit normal to a triangle
\newcommand{\Dhat}{\hat{\boldsymbol{\Delta}}}  % unit segment direction

% -------------------------------------------------------------------------
% Operations
% -------------------------------------------------------------------------
\newcommand{\norm}[1]{\left\|#1\right\|}
\newcommand{\abs}[1]{\left|#1\right|}
\newcommand{\ip}[2]{\langle #1,\, #2 \rangle}

% Projection perpendicular to a unit vector \uhat
\newcommand{\Proj}[1]{\bigl(\mathbf{I} - #1 #1^\top\bigr)}

% argmax operator (winner-take-all assignment in Phase 1 documents)
\DeclareMathOperator*{\argmax}{arg\,max}
\DeclareMathOperator*{\argmin}{arg\,min}

% -------------------------------------------------------------------------
% Calculus shorthands
% -------------------------------------------------------------------------
\newcommand{\pd}[2]{\dfrac{\partial #1}{\partial #2}}
\newcommand{\pdd}[3]{\dfrac{\partial^2 #1}{\partial #2\,\partial #3}}
\newcommand{\grad}{\nabla}
\newcommand{\Hess}[1]{\nabla^2 #1}

% -------------------------------------------------------------------------
% Optimization problem objects
% -------------------------------------------------------------------------
\newcommand{\Preg}{P_{\mathrm{reg}}}  % regular perimeter
\newcommand{\Pst}{P_{\mathrm{st}}}   % Steiner perimeter
\newcommand{\Areg}{A^{\mathrm{reg}}} % regular cell area
\newcommand{\Ast}{A^{\mathrm{st}}}   % Steiner area contribution
\newcommand{\Atgt}{\bar{A}}           % target area per cell
\newcommand{\Lagr}{\mathcal{L}}       % Lagrangian
\newcommand{\objfac}{\sigma}          % IPOPT objective scaling factor

% -------------------------------------------------------------------------
% Code cross-references (for "In the code..." callouts)
% -------------------------------------------------------------------------
\newcommand{\code}[1]{\texttt{#1}}
\newcommand{\codefile}[1]{\texttt{\small #1}}

% -------------------------------------------------------------------------
% Threshold dynamics / approach B (10-mbo-auction-dynamics)
% -------------------------------------------------------------------------
\newcommand{\Dlump}{D}                 % lumped mass matrix diag(v)
\newcommand{\Atau}{A_\tau}             % discrete diffusion operator (M + tau K)^{-1} M
\newcommand{\Etau}{E_\tau}             % heat-content (threshold-dynamics) energy
\newcommand{\Rcell}{R_{\mathrm{cell}}} % radius of the equal-area disc sqrt(Abar/pi)
\newcommand{\hmean}{h_{\mathrm{mean}}} % mean edge length
\newcommand{\hmax}{h_{\mathrm{max}}}   % max edge length
\newcommand{\Gt}{G_t}                  % heat kernel at time t
\newcommand{\Ktau}{K^{\mathrm{res}}_\tau} % resolvent (backward-Euler) kernel


% -------------------------------------------------------------------------
% Packages (include here so each document only needs % shared/macros.tex
% Shared notation for all surface-partition math documents.
% Include with: % shared/macros.tex
% Shared notation for all surface-partition math documents.
% Include with: \input{../shared/macros}

% -------------------------------------------------------------------------
% Packages (include here so each document only needs \input{../shared/macros})
% -------------------------------------------------------------------------
\usepackage{amsmath,amssymb,amsthm}
\usepackage{bm}           % \bm for bold math
\usepackage{mathtools}    % \coloneqq, dcases, etc.
\usepackage{booktabs}     % \toprule etc. in tables
\usepackage{hyperref}
\usepackage{geometry}
\usepackage{microtype}
\usepackage{parskip}
\usepackage{xcolor}
\usepackage{tcolorbox}
\usepackage{enumitem}
\usepackage{float}             % [H] float placement

\geometry{a4paper, margin=2.5cm}

% -------------------------------------------------------------------------
% Theorem-like environments
% -------------------------------------------------------------------------
\theoremstyle{definition}
\newtheorem{defn}{Definition}[section]
\newtheorem{remark}{Remark}[section]

\theoremstyle{plain}
\newtheorem{prop}{Proposition}[section]

% -------------------------------------------------------------------------
% Mesh and geometry
% -------------------------------------------------------------------------
\newcommand{\V}{\mathbf{V}}            % vertex array V[i] in R^dim
\newcommand{\F}{\mathbf{F}}            % face (triangle) array
\newcommand{\vone}[1]{v_{1,#1}}       % first vertex index of VP i
\newcommand{\vtwo}[1]{v_{2,#1}}       % second vertex index of VP i

% -------------------------------------------------------------------------
% Variable points
% -------------------------------------------------------------------------
\newcommand{\lam}{\lambda}             % scalar VP parameter
\newcommand{\Lam}{\boldsymbol{\lambda}}% full parameter vector
\newcommand{\vp}[1]{\mathbf{p}_{#1}}  % 3D position of VP i
\newcommand{\ed}[1]{\mathbf{d}_{#1}}  % edge direction d_i = V[v1_i] - V[v2_i]

% -------------------------------------------------------------------------
% Steiner / triple-point
% -------------------------------------------------------------------------
\newcommand{\Sv}{\mathbf{S}}           % Steiner (Fermat-Torricelli) point
\newcommand{\Ki}{K_{i}}               % projection matrix (I-n_i n_i^T)/r_i
\newcommand{\Mmat}{M}                 % sum of K_i matrices at a triple point

% -------------------------------------------------------------------------
% Normals and unit vectors
% -------------------------------------------------------------------------
\newcommand{\nhat}{\hat{\mathbf{n}}}  % unit normal to a triangle
\newcommand{\Dhat}{\hat{\boldsymbol{\Delta}}}  % unit segment direction

% -------------------------------------------------------------------------
% Operations
% -------------------------------------------------------------------------
\newcommand{\norm}[1]{\left\|#1\right\|}
\newcommand{\abs}[1]{\left|#1\right|}
\newcommand{\ip}[2]{\langle #1,\, #2 \rangle}

% Projection perpendicular to a unit vector \uhat
\newcommand{\Proj}[1]{\bigl(\mathbf{I} - #1 #1^\top\bigr)}

% argmax operator (winner-take-all assignment in Phase 1 documents)
\DeclareMathOperator*{\argmax}{arg\,max}
\DeclareMathOperator*{\argmin}{arg\,min}

% -------------------------------------------------------------------------
% Calculus shorthands
% -------------------------------------------------------------------------
\newcommand{\pd}[2]{\dfrac{\partial #1}{\partial #2}}
\newcommand{\pdd}[3]{\dfrac{\partial^2 #1}{\partial #2\,\partial #3}}
\newcommand{\grad}{\nabla}
\newcommand{\Hess}[1]{\nabla^2 #1}

% -------------------------------------------------------------------------
% Optimization problem objects
% -------------------------------------------------------------------------
\newcommand{\Preg}{P_{\mathrm{reg}}}  % regular perimeter
\newcommand{\Pst}{P_{\mathrm{st}}}   % Steiner perimeter
\newcommand{\Areg}{A^{\mathrm{reg}}} % regular cell area
\newcommand{\Ast}{A^{\mathrm{st}}}   % Steiner area contribution
\newcommand{\Atgt}{\bar{A}}           % target area per cell
\newcommand{\Lagr}{\mathcal{L}}       % Lagrangian
\newcommand{\objfac}{\sigma}          % IPOPT objective scaling factor

% -------------------------------------------------------------------------
% Code cross-references (for "In the code..." callouts)
% -------------------------------------------------------------------------
\newcommand{\code}[1]{\texttt{#1}}
\newcommand{\codefile}[1]{\texttt{\small #1}}

% -------------------------------------------------------------------------
% Threshold dynamics / approach B (10-mbo-auction-dynamics)
% -------------------------------------------------------------------------
\newcommand{\Dlump}{D}                 % lumped mass matrix diag(v)
\newcommand{\Atau}{A_\tau}             % discrete diffusion operator (M + tau K)^{-1} M
\newcommand{\Etau}{E_\tau}             % heat-content (threshold-dynamics) energy
\newcommand{\Rcell}{R_{\mathrm{cell}}} % radius of the equal-area disc sqrt(Abar/pi)
\newcommand{\hmean}{h_{\mathrm{mean}}} % mean edge length
\newcommand{\hmax}{h_{\mathrm{max}}}   % max edge length
\newcommand{\Gt}{G_t}                  % heat kernel at time t
\newcommand{\Ktau}{K^{\mathrm{res}}_\tau} % resolvent (backward-Euler) kernel


% -------------------------------------------------------------------------
% Packages (include here so each document only needs % shared/macros.tex
% Shared notation for all surface-partition math documents.
% Include with: \input{../shared/macros}

% -------------------------------------------------------------------------
% Packages (include here so each document only needs \input{../shared/macros})
% -------------------------------------------------------------------------
\usepackage{amsmath,amssymb,amsthm}
\usepackage{bm}           % \bm for bold math
\usepackage{mathtools}    % \coloneqq, dcases, etc.
\usepackage{booktabs}     % \toprule etc. in tables
\usepackage{hyperref}
\usepackage{geometry}
\usepackage{microtype}
\usepackage{parskip}
\usepackage{xcolor}
\usepackage{tcolorbox}
\usepackage{enumitem}
\usepackage{float}             % [H] float placement

\geometry{a4paper, margin=2.5cm}

% -------------------------------------------------------------------------
% Theorem-like environments
% -------------------------------------------------------------------------
\theoremstyle{definition}
\newtheorem{defn}{Definition}[section]
\newtheorem{remark}{Remark}[section]

\theoremstyle{plain}
\newtheorem{prop}{Proposition}[section]

% -------------------------------------------------------------------------
% Mesh and geometry
% -------------------------------------------------------------------------
\newcommand{\V}{\mathbf{V}}            % vertex array V[i] in R^dim
\newcommand{\F}{\mathbf{F}}            % face (triangle) array
\newcommand{\vone}[1]{v_{1,#1}}       % first vertex index of VP i
\newcommand{\vtwo}[1]{v_{2,#1}}       % second vertex index of VP i

% -------------------------------------------------------------------------
% Variable points
% -------------------------------------------------------------------------
\newcommand{\lam}{\lambda}             % scalar VP parameter
\newcommand{\Lam}{\boldsymbol{\lambda}}% full parameter vector
\newcommand{\vp}[1]{\mathbf{p}_{#1}}  % 3D position of VP i
\newcommand{\ed}[1]{\mathbf{d}_{#1}}  % edge direction d_i = V[v1_i] - V[v2_i]

% -------------------------------------------------------------------------
% Steiner / triple-point
% -------------------------------------------------------------------------
\newcommand{\Sv}{\mathbf{S}}           % Steiner (Fermat-Torricelli) point
\newcommand{\Ki}{K_{i}}               % projection matrix (I-n_i n_i^T)/r_i
\newcommand{\Mmat}{M}                 % sum of K_i matrices at a triple point

% -------------------------------------------------------------------------
% Normals and unit vectors
% -------------------------------------------------------------------------
\newcommand{\nhat}{\hat{\mathbf{n}}}  % unit normal to a triangle
\newcommand{\Dhat}{\hat{\boldsymbol{\Delta}}}  % unit segment direction

% -------------------------------------------------------------------------
% Operations
% -------------------------------------------------------------------------
\newcommand{\norm}[1]{\left\|#1\right\|}
\newcommand{\abs}[1]{\left|#1\right|}
\newcommand{\ip}[2]{\langle #1,\, #2 \rangle}

% Projection perpendicular to a unit vector \uhat
\newcommand{\Proj}[1]{\bigl(\mathbf{I} - #1 #1^\top\bigr)}

% argmax operator (winner-take-all assignment in Phase 1 documents)
\DeclareMathOperator*{\argmax}{arg\,max}
\DeclareMathOperator*{\argmin}{arg\,min}

% -------------------------------------------------------------------------
% Calculus shorthands
% -------------------------------------------------------------------------
\newcommand{\pd}[2]{\dfrac{\partial #1}{\partial #2}}
\newcommand{\pdd}[3]{\dfrac{\partial^2 #1}{\partial #2\,\partial #3}}
\newcommand{\grad}{\nabla}
\newcommand{\Hess}[1]{\nabla^2 #1}

% -------------------------------------------------------------------------
% Optimization problem objects
% -------------------------------------------------------------------------
\newcommand{\Preg}{P_{\mathrm{reg}}}  % regular perimeter
\newcommand{\Pst}{P_{\mathrm{st}}}   % Steiner perimeter
\newcommand{\Areg}{A^{\mathrm{reg}}} % regular cell area
\newcommand{\Ast}{A^{\mathrm{st}}}   % Steiner area contribution
\newcommand{\Atgt}{\bar{A}}           % target area per cell
\newcommand{\Lagr}{\mathcal{L}}       % Lagrangian
\newcommand{\objfac}{\sigma}          % IPOPT objective scaling factor

% -------------------------------------------------------------------------
% Code cross-references (for "In the code..." callouts)
% -------------------------------------------------------------------------
\newcommand{\code}[1]{\texttt{#1}}
\newcommand{\codefile}[1]{\texttt{\small #1}}

% -------------------------------------------------------------------------
% Threshold dynamics / approach B (10-mbo-auction-dynamics)
% -------------------------------------------------------------------------
\newcommand{\Dlump}{D}                 % lumped mass matrix diag(v)
\newcommand{\Atau}{A_\tau}             % discrete diffusion operator (M + tau K)^{-1} M
\newcommand{\Etau}{E_\tau}             % heat-content (threshold-dynamics) energy
\newcommand{\Rcell}{R_{\mathrm{cell}}} % radius of the equal-area disc sqrt(Abar/pi)
\newcommand{\hmean}{h_{\mathrm{mean}}} % mean edge length
\newcommand{\hmax}{h_{\mathrm{max}}}   % max edge length
\newcommand{\Gt}{G_t}                  % heat kernel at time t
\newcommand{\Ktau}{K^{\mathrm{res}}_\tau} % resolvent (backward-Euler) kernel
)
% -------------------------------------------------------------------------
\usepackage{amsmath,amssymb,amsthm}
\usepackage{bm}           % \bm for bold math
\usepackage{mathtools}    % \coloneqq, dcases, etc.
\usepackage{booktabs}     % \toprule etc. in tables
\usepackage{hyperref}
\usepackage{geometry}
\usepackage{microtype}
\usepackage{parskip}
\usepackage{xcolor}
\usepackage{tcolorbox}
\usepackage{enumitem}
\usepackage{float}             % [H] float placement

\geometry{a4paper, margin=2.5cm}

% -------------------------------------------------------------------------
% Theorem-like environments
% -------------------------------------------------------------------------
\theoremstyle{definition}
\newtheorem{defn}{Definition}[section]
\newtheorem{remark}{Remark}[section]

\theoremstyle{plain}
\newtheorem{prop}{Proposition}[section]

% -------------------------------------------------------------------------
% Mesh and geometry
% -------------------------------------------------------------------------
\newcommand{\V}{\mathbf{V}}            % vertex array V[i] in R^dim
\newcommand{\F}{\mathbf{F}}            % face (triangle) array
\newcommand{\vone}[1]{v_{1,#1}}       % first vertex index of VP i
\newcommand{\vtwo}[1]{v_{2,#1}}       % second vertex index of VP i

% -------------------------------------------------------------------------
% Variable points
% -------------------------------------------------------------------------
\newcommand{\lam}{\lambda}             % scalar VP parameter
\newcommand{\Lam}{\boldsymbol{\lambda}}% full parameter vector
\newcommand{\vp}[1]{\mathbf{p}_{#1}}  % 3D position of VP i
\newcommand{\ed}[1]{\mathbf{d}_{#1}}  % edge direction d_i = V[v1_i] - V[v2_i]

% -------------------------------------------------------------------------
% Steiner / triple-point
% -------------------------------------------------------------------------
\newcommand{\Sv}{\mathbf{S}}           % Steiner (Fermat-Torricelli) point
\newcommand{\Ki}{K_{i}}               % projection matrix (I-n_i n_i^T)/r_i
\newcommand{\Mmat}{M}                 % sum of K_i matrices at a triple point

% -------------------------------------------------------------------------
% Normals and unit vectors
% -------------------------------------------------------------------------
\newcommand{\nhat}{\hat{\mathbf{n}}}  % unit normal to a triangle
\newcommand{\Dhat}{\hat{\boldsymbol{\Delta}}}  % unit segment direction

% -------------------------------------------------------------------------
% Operations
% -------------------------------------------------------------------------
\newcommand{\norm}[1]{\left\|#1\right\|}
\newcommand{\abs}[1]{\left|#1\right|}
\newcommand{\ip}[2]{\langle #1,\, #2 \rangle}

% Projection perpendicular to a unit vector \uhat
\newcommand{\Proj}[1]{\bigl(\mathbf{I} - #1 #1^\top\bigr)}

% argmax operator (winner-take-all assignment in Phase 1 documents)
\DeclareMathOperator*{\argmax}{arg\,max}
\DeclareMathOperator*{\argmin}{arg\,min}

% -------------------------------------------------------------------------
% Calculus shorthands
% -------------------------------------------------------------------------
\newcommand{\pd}[2]{\dfrac{\partial #1}{\partial #2}}
\newcommand{\pdd}[3]{\dfrac{\partial^2 #1}{\partial #2\,\partial #3}}
\newcommand{\grad}{\nabla}
\newcommand{\Hess}[1]{\nabla^2 #1}

% -------------------------------------------------------------------------
% Optimization problem objects
% -------------------------------------------------------------------------
\newcommand{\Preg}{P_{\mathrm{reg}}}  % regular perimeter
\newcommand{\Pst}{P_{\mathrm{st}}}   % Steiner perimeter
\newcommand{\Areg}{A^{\mathrm{reg}}} % regular cell area
\newcommand{\Ast}{A^{\mathrm{st}}}   % Steiner area contribution
\newcommand{\Atgt}{\bar{A}}           % target area per cell
\newcommand{\Lagr}{\mathcal{L}}       % Lagrangian
\newcommand{\objfac}{\sigma}          % IPOPT objective scaling factor

% -------------------------------------------------------------------------
% Code cross-references (for "In the code..." callouts)
% -------------------------------------------------------------------------
\newcommand{\code}[1]{\texttt{#1}}
\newcommand{\codefile}[1]{\texttt{\small #1}}

% -------------------------------------------------------------------------
% Threshold dynamics / approach B (10-mbo-auction-dynamics)
% -------------------------------------------------------------------------
\newcommand{\Dlump}{D}                 % lumped mass matrix diag(v)
\newcommand{\Atau}{A_\tau}             % discrete diffusion operator (M + tau K)^{-1} M
\newcommand{\Etau}{E_\tau}             % heat-content (threshold-dynamics) energy
\newcommand{\Rcell}{R_{\mathrm{cell}}} % radius of the equal-area disc sqrt(Abar/pi)
\newcommand{\hmean}{h_{\mathrm{mean}}} % mean edge length
\newcommand{\hmax}{h_{\mathrm{max}}}   % max edge length
\newcommand{\Gt}{G_t}                  % heat kernel at time t
\newcommand{\Ktau}{K^{\mathrm{res}}_\tau} % resolvent (backward-Euler) kernel
)
% -------------------------------------------------------------------------
\usepackage{amsmath,amssymb,amsthm}
\usepackage{bm}           % \bm for bold math
\usepackage{mathtools}    % \coloneqq, dcases, etc.
\usepackage{booktabs}     % \toprule etc. in tables
\usepackage{hyperref}
\usepackage{geometry}
\usepackage{microtype}
\usepackage{parskip}
\usepackage{xcolor}
\usepackage{tcolorbox}
\usepackage{enumitem}
\usepackage{float}             % [H] float placement

\geometry{a4paper, margin=2.5cm}

% -------------------------------------------------------------------------
% Theorem-like environments
% -------------------------------------------------------------------------
\theoremstyle{definition}
\newtheorem{defn}{Definition}[section]
\newtheorem{remark}{Remark}[section]

\theoremstyle{plain}
\newtheorem{prop}{Proposition}[section]

% -------------------------------------------------------------------------
% Mesh and geometry
% -------------------------------------------------------------------------
\newcommand{\V}{\mathbf{V}}            % vertex array V[i] in R^dim
\newcommand{\F}{\mathbf{F}}            % face (triangle) array
\newcommand{\vone}[1]{v_{1,#1}}       % first vertex index of VP i
\newcommand{\vtwo}[1]{v_{2,#1}}       % second vertex index of VP i

% -------------------------------------------------------------------------
% Variable points
% -------------------------------------------------------------------------
\newcommand{\lam}{\lambda}             % scalar VP parameter
\newcommand{\Lam}{\boldsymbol{\lambda}}% full parameter vector
\newcommand{\vp}[1]{\mathbf{p}_{#1}}  % 3D position of VP i
\newcommand{\ed}[1]{\mathbf{d}_{#1}}  % edge direction d_i = V[v1_i] - V[v2_i]

% -------------------------------------------------------------------------
% Steiner / triple-point
% -------------------------------------------------------------------------
\newcommand{\Sv}{\mathbf{S}}           % Steiner (Fermat-Torricelli) point
\newcommand{\Ki}{K_{i}}               % projection matrix (I-n_i n_i^T)/r_i
\newcommand{\Mmat}{M}                 % sum of K_i matrices at a triple point

% -------------------------------------------------------------------------
% Normals and unit vectors
% -------------------------------------------------------------------------
\newcommand{\nhat}{\hat{\mathbf{n}}}  % unit normal to a triangle
\newcommand{\Dhat}{\hat{\boldsymbol{\Delta}}}  % unit segment direction

% -------------------------------------------------------------------------
% Operations
% -------------------------------------------------------------------------
\newcommand{\norm}[1]{\left\|#1\right\|}
\newcommand{\abs}[1]{\left|#1\right|}
\newcommand{\ip}[2]{\langle #1,\, #2 \rangle}

% Projection perpendicular to a unit vector \uhat
\newcommand{\Proj}[1]{\bigl(\mathbf{I} - #1 #1^\top\bigr)}

% argmax operator (winner-take-all assignment in Phase 1 documents)
\DeclareMathOperator*{\argmax}{arg\,max}
\DeclareMathOperator*{\argmin}{arg\,min}

% -------------------------------------------------------------------------
% Calculus shorthands
% -------------------------------------------------------------------------
\newcommand{\pd}[2]{\dfrac{\partial #1}{\partial #2}}
\newcommand{\pdd}[3]{\dfrac{\partial^2 #1}{\partial #2\,\partial #3}}
\newcommand{\grad}{\nabla}
\newcommand{\Hess}[1]{\nabla^2 #1}

% -------------------------------------------------------------------------
% Optimization problem objects
% -------------------------------------------------------------------------
\newcommand{\Preg}{P_{\mathrm{reg}}}  % regular perimeter
\newcommand{\Pst}{P_{\mathrm{st}}}   % Steiner perimeter
\newcommand{\Areg}{A^{\mathrm{reg}}} % regular cell area
\newcommand{\Ast}{A^{\mathrm{st}}}   % Steiner area contribution
\newcommand{\Atgt}{\bar{A}}           % target area per cell
\newcommand{\Lagr}{\mathcal{L}}       % Lagrangian
\newcommand{\objfac}{\sigma}          % IPOPT objective scaling factor

% -------------------------------------------------------------------------
% Code cross-references (for "In the code..." callouts)
% -------------------------------------------------------------------------
\newcommand{\code}[1]{\texttt{#1}}
\newcommand{\codefile}[1]{\texttt{\small #1}}

% -------------------------------------------------------------------------
% Threshold dynamics / approach B (10-mbo-auction-dynamics)
% -------------------------------------------------------------------------
\newcommand{\Dlump}{D}                 % lumped mass matrix diag(v)
\newcommand{\Atau}{A_\tau}             % discrete diffusion operator (M + tau K)^{-1} M
\newcommand{\Etau}{E_\tau}             % heat-content (threshold-dynamics) energy
\newcommand{\Rcell}{R_{\mathrm{cell}}} % radius of the equal-area disc sqrt(Abar/pi)
\newcommand{\hmean}{h_{\mathrm{mean}}} % mean edge length
\newcommand{\hmax}{h_{\mathrm{max}}}   % max edge length
\newcommand{\Gt}{G_t}                  % heat kernel at time t
\newcommand{\Ktau}{K^{\mathrm{res}}_\tau} % resolvent (backward-Euler) kernel


\hypersetup{
  colorlinks = true,
  linkcolor  = blue!60!black,
  citecolor  = green!50!black,
  urlcolor   = blue!60!black,
  pdftitle   = {Phase 1 Scaling with Number of Regions},
  pdfauthor  = {surface-partition project}
}

\title{\textbf{Phase~1 Wall-Time Scaling with Number of Regions}\\[0.5em]
  \large Empirical Measurements, a Corrected Power-Law Fit, and the Path to $N>100$}
\author{}
\date{July 2026}

\begin{document}
\maketitle

\begin{abstract}
This document analyses how Phase~1 (PGD relaxation) wall time grows with the
number of regions~$N$, using profiled runs on the torus.  An earlier draft of
this document fit a power law of exponent $\alpha \approx 3.1$ ($T \sim N^3$)
to two \emph{pre-optimization} points --- $N=10$ (random init, 7 levels) and
$N=30$ (seeded init, 5 levels) --- and extrapolated it to conclude that
$N=1000$ would take $\sim\!14$ years and is ``not achievable.''  \textbf{That
conclusion was an artifact of the fit and is retracted here.}  The two points
confounded three variables at once (random vs.\ seeded initialisation, $7$
vs.\ $5$ refinement levels, and pre- vs.\ post-optimization code), and the
accompanying claim that the outer iteration count $K$ grows as $\sim\!N^2$ is
contradicted by direct measurement.

New post-optimization, same-init (seeded) runs at $N=50$ and $N=100$ pin the
scaling far more honestly.  At an equal finest mesh
($348\times328$, $V=114{,}144$), $K$ is essentially \emph{flat} ---
$54{,}154$ at $N=50$ versus $54{,}536$ at $N=100$ --- not $\sim\!N^2$ (which
would predict a $4\times$ increase).  Total wall grows from
$\mathbf{11.6}$~h ($N=50$) to $\mathbf{40.6}$~h ($N=100$) at equal mesh, an
effective exponent $\alpha\approx1.8$; matched instead at constant per-cell
resolution ($\sim\!2290$ vertices/cell) it grows from $11.6$~h to
$\mathbf{29.3}$~h, $\alpha\approx1.3$ --- but that finer run also converged in
about a third of the outer iterations, so $1.3$ is a cost-at-fixed-quality
figure that \emph{understates} the intrinsic $N$-exponent, which the cost model
places nearer $\alpha\approx1.8$--$2.2$ (\S\ref{sec:fit}).  Both $N=100$ runs
produce \emph{valid} $100$-region partitions.  What actually grows with $N$ is the
per-call projection cost --- the mean projection inner-iteration count
$\mathrm{mPII}$ climbs $13.8\to21\to33$ --- while the outer iteration count
stays flat.

The corrected outlook: post-optimization seeded scaling between $N=50$ and
$N=100$ measures $\alpha\approx1.3$--$1.8$, and the intrinsic (fixed-iteration)
$N$-exponent likely sits at or above the top of that band ($\approx1.8$--$2.2$).
Under $\alpha\approx1.5$--$2$ this puts $N=1000$ at a matter of
\textbf{weeks-to-months} (an illustrative $\alpha\approx1.5$ gives $\sim\!43$
days; $\alpha\approx2$ gives $\sim\!190$ days), \emph{not} years.  That figure is
an extrapolation from only two-to-three post-optimization points over a
$20\times$ range in~$N$ and cannot be trusted beyond an order of magnitude; the
observed $\mathrm{mPII}$ growth could steepen the effective exponent further at
large~$N$.  The essential correction is
that the catastrophic ``$\sim\!14$ years / not achievable'' framing is
\textbf{not supported by the data}.  The dominant cost remains the orthogonal
projection ($80$--$87\%$ of wall), whose per-call cost scales as
$\mathcal{O}(\mathrm{mPII}\cdot VN)$; the durable algorithmic fix is a
closed-form simplex projection with a dual-Newton equal-area coupling
(\S\ref{sec:improvements}), which attacks $\mathrm{mPII}$ directly.
\end{abstract}

\tableofcontents
\newpage

% =========================================================================
\section{Context and Motivation}
% =========================================================================

The long-term goal of the surface-partition project is to compute
minimal-perimeter partitions for $N \gg 30$ regions, with $N = 1000$ as a
target.  Document~04 (\emph{Phase~1 Timing Profile}) established that
$80$--$90\%$ of Phase~1 wall time is spent in the orthogonal projection
$\Pi_{\mathcal{C}}$ onto the partition-of-unity and equal-area constraints,
and that the projection cost per call is $\mathcal{O}(\mathrm{mPII} \cdot V
N)$, where $V$ is the number of mesh vertices and $\mathrm{mPII}$ is the
mean inner-iteration count.

This document quantifies how total Phase~1 wall time actually grows with $N$
using the $N$ values for which profiled data exist, identifies which factors
in the cost model are responsible for the growth, and projects --- with
explicit caveats --- to larger~$N$.

\begin{remark}[A retracted conclusion]
\label{rem:retraction}
An earlier version of this document concluded $\alpha\approx3.1$ ($T\sim N^3$)
and, from it, $N=1000\approx14$~years and ``not achievable.''  Both figures are
\textbf{withdrawn}.  They came from a two-point fit
(\S\ref{sec:why-old-fit-wrong}) that mixed initialisation methods, level
counts, and pre- vs.\ post-optimization code, and that assumed an empirically
false $K\sim N^2$ iteration-count growth.  The corrected analysis below uses
only the post-optimization, same-init runs at $N=50$ and $N=100$, which are the
first genuinely comparable pair.
\end{remark}

\begin{remark}[Optimization status of each run]
The $N=10$ and $N=30$ runs of \S\ref{sec:data} predate the Phase~1 PGD serial
optimizations (Changes~A/B/C; see \S\ref{sec:serial-opts} and
\texttt{docs/reference/PHASE1\_PGD\_SERIAL\_OPTIMIZATION\_AUDIT.md}).  The
$N=50$ run (\S\ref{sec:n50}) and \emph{both} $N=100$ runs (\S\ref{sec:n100})
are post-optimization.  Only the post-optimization, same-init points are used
to estimate the exponent; the pre-optimization points are retained for context
only and are explicitly \emph{not} fit together with the post-optimization
data.
\end{remark}

% =========================================================================
\section{Notation and Setup}
\label{sec:notation}
% =========================================================================

Throughout this document:

\begin{itemize}
  \item $N$ --- number of regions (partition cells).
  \item $V$ --- number of mesh vertices at the finest refinement level.
  \item $L$ --- number of refinement levels.
  \item $T$ --- total Phase~1 wall time (seconds or hours).  Taken from
    \texttt{summary.total\_wall\_s} in \texttt{solution/timing\_profile.yaml}.
  \item $\mathrm{mBT}$ --- mean Armijo backtrack trial steps per major iteration.
  \item $\mathrm{mPII}$ --- mean inner-iteration count per projection call.
  \item $K$ --- total major (outer) PGD iterations summed across all levels.
  \item ``V/cell'' --- finest-level vertices per region, $V/N$, a proxy for the
    per-cell mesh resolution.
\end{itemize}

The cost model from document~04 is
\begin{equation}
  T \;\sim\;
  \underbrace{K \cdot (\mathrm{mBT}+1)}_{\text{\# projection calls}}
  \;\times\;
  \underbrace{\mathrm{mPII} \cdot V N}_{\text{cost per call}}.
  \label{eq:cost-model}
\end{equation}
The analysis below tracks each factor separately across $N$, rather than
collapsing them into a single exponent up front.

% =========================================================================
\section{Measured Data Points}
\label{sec:data}
% =========================================================================

All wall times are read from \texttt{summary.total\_wall\_s}; all validity
statements are read from the \texttt{dormant\_cells} block of
\texttt{solution/metadata.yaml} (a partition is valid when both \texttt{dead}
and \texttt{weak} lists are empty).

\subsection{$N = 10$ and $N=30$: pre-optimization context (not fittable)}
\label{sec:pre-opt}

Two families of pre-optimization runs exist.  They are recorded here for
context, but --- being pre-optimization and, in the $N=10$ case, differently
initialised --- they \textbf{must not} be fit together with the
post-optimization points to produce an exponent (see
\S\ref{sec:why-old-fit-wrong}).

The $N=10$ reference run used the \emph{random} initialisation, 7 refinement
levels, and an $80\times66$ initial mesh growing to $V_{10} = 187{,}128$
vertices; total wall time $T_{10} = 1.19$\,h.  Per-level iteration counts were
not recorded (older schema).

The $N=30$ family used the \emph{seeded} (Voronoi farthest-point)
initialisation, 5 levels, and a $90\times86$ initial mesh growing to
$V_{30} = 107{,}484$ vertices.  All produced valid 30-region partitions.  The
best-case run is seed 84172851: $T_{30}=74{,}670.9\,\text{s}=\mathbf{20.74}$\,h,
$K=27{,}001$, $\mathrm{mBT}=6.26$, $\mathrm{mPII}=19.39$, projection share
$89.4\%$, valid (\texttt{dead}$=[\,]$, \texttt{weak}$=[\,]$).

\begin{table}[H]
\centering
\caption{Pre-optimization reference runs (torus).  Retained for context only;
\emph{not} used to estimate the exponent.  The $N=10$ run additionally differs
in initialisation (random) and level count (7).}
\label{tab:pre-opt}
\renewcommand{\arraystretch}{1.2}
\begin{tabular}{rllrrl}
\toprule
$N$ & init & levels & $V$ (finest) & Wall (h) & valid \\
\midrule
10 & random & 7 & 187\,128 & \phantom{0}1.19 & yes \\
30 & seeded & 5 & 107\,484 & 20.74 & yes \\
\bottomrule
\end{tabular}
\end{table}

\subsection{$N = 50$ measured run (post-optimization, seeded)}
\label{sec:n50}

A profiled run taken \emph{after} the serial optimizations of
\S\ref{sec:serial-opts} provides the first post-optimization anchor.  It used
the seeded initialisation, 5 refinement levels, and a $100\times96$ initial
mesh growing to $V_{50}=114{,}144$ vertices ($2{,}283$ V/cell).  On a clean
host (seed 84172851) the total wall time was
$T_{50}=41{,}754.3\,\text{s} = \mathbf{11.6}$\,h.  The partition is valid
(\texttt{dead}$=[\,]$, \texttt{weak}$=[\,]$; minimum peak density $\approx1.0$).

\begin{table}[H]
\centering
\caption{$N=50$ post-optimization run (torus, $\lambda=2.1$, seed 84172851,
5 levels, seeded init, $V_{50}=114{,}144$ vertices; run
\texttt{run\_20260627\_234511}).}
\label{tab:n50}
\renewcommand{\arraystretch}{1.2}
\begin{tabular}{ll}
\toprule
\textbf{Quantity} & \textbf{Value} \\
\midrule
Total wall time $T_{50}$       & 11.6\,h ($41{,}754$\,s) \\
Major iterations $K$           & 54\,154 \\
Mean backtracks $\mathrm{mBT}$ & 1.93 \\
Mean inner iters $\mathrm{mPII}$ & 13.76 \\
Projection share               & 80.5\,\% of wall \\
Energy / gradient share        & 8.2\,\% / 7.4\,\% \\
Backtrack share (net)          & 2.5\,\% \\
Validity                       & valid ($\mathrm{dead}=\mathrm{weak}=\varnothing$) \\
\bottomrule
\end{tabular}
\end{table}

\subsection{$N = 100$ measured runs (post-optimization, seeded)}
\label{sec:n100}

Two post-optimization $N=100$ runs, both seeded and both producing
\emph{valid} 100-region partitions, form the first genuinely comparable pair
with $N=50$:

\begin{itemize}
  \item \textbf{Equal-mesh run} (\texttt{run\_20260629\_141012}): the
    \emph{same} $348\times328$ finest mesh as the $N=50$ run
    ($V=114{,}144$, hence $1{,}141$ V/cell --- under-resolved for $100$
    cells but still valid, minimum peak density $0.991$).  This isolates the
    $N$-dependence of the algorithm at fixed mesh.
  \item \textbf{Resolution-matched run} (\texttt{run\_20260701\_143238}): a
    finer $494\times464$ mesh, $V=229{,}216$, $2{,}292$ V/cell --- matched to
    the $N=50$ run's per-cell resolution.  Valid, minimum peak density
    $0.976$.  This is the fair ``same partition quality'' comparison, but it
    also doubles~$V$.
\end{itemize}

\begin{table}[H]
\centering
\caption{Post-optimization $N=100$ runs (torus, $\lambda=2.1$, seed 84172851,
5 levels, seeded init).  Both are valid partitions
($\mathrm{dead}=\mathrm{weak}=\varnothing$).  The $N=50$ row is repeated for
comparison.  ``$K$ flat'' is the key observation: at equal mesh $K$ barely
moves.}
\label{tab:n100}
\renewcommand{\arraystretch}{1.25}
\begin{tabular}{lrrrrr}
\toprule
Run & $N$ & finest mesh & $V$ & V/cell & Wall (h) \\
\midrule
\texttt{234511} & \phantom{0}50 & $348\times328$ & 114\,144 & 2\,283 & 11.6 \\
\texttt{141012} & 100 & $348\times328$ & 114\,144 & 1\,141 & 40.6 \\
\texttt{143238} & 100 & $494\times464$ & 229\,216 & 2\,292 & 29.3 \\
\bottomrule
\end{tabular}
\end{table}

\begin{table}[H]
\centering
\caption{Profile detail for the same three runs.  Note $K$ (major iterations)
is flat between the two equal-mesh runs while $\mathrm{mPII}$ (per-call
projection inner iterations) climbs steadily with $N$ and mesh size.}
\label{tab:n100-profile}
\renewcommand{\arraystretch}{1.25}
\begin{tabular}{lrrrrr}
\toprule
Run & $N$ / mesh & $K$ & $\mathrm{mBT}$ & $\mathrm{mPII}$ & proj \% \\
\midrule
\texttt{234511} & 50 / $348^2$-class  & 54\,154 & 1.93 & 13.76 & 80.5 \\
\texttt{141012} & 100 / same mesh     & 54\,536 & 1.96 & 21.15 & 86.0 \\
\texttt{143238} & 100 / finer mesh    & 16\,908 & 1.48 & 32.95 & 87.1 \\
\bottomrule
\end{tabular}
\end{table}

The single most important number in this document is in
Table~\ref{tab:n100-profile}: at an \emph{equal} finest mesh, doubling $N$ from
$50$ to $100$ moved $K$ from $54{,}154$ to $54{,}536$ --- a $0.7\%$ change.  The
old $K\sim N^2$ assumption would have predicted $K$ near $216{,}000$
(a $4\times$ increase).  The iteration count does \emph{not} blow up with $N$.

% =========================================================================
\section{Why the Old $N^3$ / 14-Year Fit Was an Artifact}
\label{sec:why-old-fit-wrong}
% =========================================================================

The retracted $\alpha\approx3.1$ came from fitting $T\propto N^{\alpha}$ to
exactly two points, $N=10$ and $N=30$, after mesh-normalising $T_{10}$.  Three
independent defects make that fit unsafe, and a fourth defect propagated into
the projection.

\paragraph{Defect 1 --- only two points.}
A power law has two free parameters; fitting it to two points leaves zero
degrees of freedom, so the exponent is whatever the pair happens to imply, with
no residual to expose a bad fit.

\paragraph{Defect 2 --- the initialisation confound.}
$N=10$ used the \emph{random} initialisation; $N=30$ used the \emph{seeded}
initialisation.  These are different algorithms with different convergence
behaviour (the seeded path hands each cell a contiguous winning region from
iteration~0).  The observed $T_{10}\to T_{30}$ jump therefore blends an
$N$-effect with an init-method effect, and the two cannot be separated from a
single pair.

\paragraph{Defect 3 --- the pre- vs.\ post-optimization confound, and level
counts.}
Both points predate Changes~A/B/C, and they used different level counts (7 vs.\
5).  The serial optimizations changed $\mathrm{mBT}$ and the projection share
substantially (\S\ref{sec:serial-opts}); a scaling law calibrated on
pre-optimization runs does not transfer cleanly to the current code.

\paragraph{Defect 4 --- the empirically false $K\sim N^2$ assumption.}
The old derivation decomposed $\alpha\approx3$ as ``$\mathcal{O}(N)$ per call
$\times$ $\mathcal{O}(N^2)$ iteration count.''  The $N^2$ factor was
\emph{inferred from the same confounded pair} (an estimated $K\approx3{,}500$ at
$N=10$ against a measured $K\approx27{,}000$--$50{,}000$ at $N=30$), not
measured at fixed init and code.  The post-optimization same-init data refute
it directly: at equal mesh, $K = 54{,}154 \to 54{,}536$ from $N=50$ to $N=100$
(Table~\ref{tab:n100-profile}).  $K$ is flat, so the $N^2$ term --- the sole
source of the ``extra'' two powers of $N$ that produced $N^3$ --- does not
exist in the measured regime.

\medskip
With the $K\sim N^2$ factor removed, the surviving $N$-dependence is the
per-call cost of the projection, which the cost model
\eqref{eq:cost-model} puts at $\mathcal{O}(\mathrm{mPII}\cdot VN)$.  That is a
low-order growth, not a cubic one, and it is what the next section measures.

% =========================================================================
\section{Corrected Empirical Scaling}
\label{sec:fit}
% =========================================================================

We estimate the exponent using \emph{only} the post-optimization, same-init
(seeded) points: the $N=50$ run and the two $N=100$ runs.  Two honest
comparisons are available, and they measure different things.

\subsection{Equal-mesh comparison (fixed $V$)}

Holding the finest mesh fixed at $348\times328$ ($V=114{,}144$) and doubling
$N$ from $50$ to $100$:
\begin{equation}
  \frac{T_{100}^{\text{eq}}}{T_{50}}
  = \frac{146{,}240.4\,\text{s}}{41{,}754.3\,\text{s}}
  = 3.50,
  \qquad
  \alpha_{\text{eq}}
  = \frac{\ln 3.50}{\ln 2}
  = \frac{1.253}{0.693}
  \approx \mathbf{1.81}.
  \label{eq:alpha-eq}
\end{equation}
This isolates the pure $N$-effect at fixed mesh work.  It is an \emph{upper}
estimate of the practically relevant exponent, because at fixed $V$ the
$N=100$ partition is under-resolved ($1{,}141$ V/cell), which forces more
projection inner iterations per call ($\mathrm{mPII}$: $13.76\to21.15$).

\subsection{Resolution-matched comparison (fixed V/cell)}

Comparing instead at constant per-cell resolution --- the $N=50$ run
($V=114{,}144$, $2{,}283$ V/cell) against the finer $N=100$ run
($V=229{,}216$, $2{,}292$ V/cell):
\begin{equation}
  \frac{T_{100}^{\text{fine}}}{T_{50}}
  = \frac{105{,}601.0\,\text{s}}{41{,}754.3\,\text{s}}
  = 2.53,
  \qquad
  \alpha_{\text{rm}}
  = \frac{\ln 2.53}{\ln 2}
  = \frac{0.928}{0.693}
  \approx \mathbf{1.34}.
  \label{eq:alpha-rm}
\end{equation}
This is the fair ``same partition quality'' comparison --- adding regions in
practice requires a proportionally finer mesh, and it folds in the mesh growth
($V$ also doubled here).  But $\alpha_{\text{rm}}=1.34$ must \textbf{not} be read
as a clean pure-$N$ exponent: the finer run is fast largely because its outer
iteration count \emph{collapsed} to $K=16{,}908$, about a third of the
$\sim\!54{,}000$ of the other two runs (Table~\ref{tab:n100-profile}), with
fewer backtracks ($\mathrm{mBT}$ $1.48$ vs.\ $1.93$).  That $K$ drop is
unexplained and may be specific to this mesh/level schedule rather than an
$N$-scaling property.  Crucially it \emph{contradicts} the ``$K$ flat''
finding this document leans on elsewhere: had $K$ stayed at $54{,}154$, the
resolution-matched ratio would be $8.1$ and $\alpha_{\text{rm}}\approx3.0$.
Indeed, holding $K$ fixed with $V\propto N$ (constant V/cell), the cost model
\eqref{eq:cost-model} gives $T\sim\mathrm{mPII}\cdot V N\sim\mathrm{mPII}\cdot
N^2$, i.e.\ $\alpha_{\text{rm}}\gtrsim2$ before the $\mathrm{mPII}$ rise is even
counted.  So $1.34$ is a \emph{lower engineering bound} (cost at fixed quality,
this run's convergence included), while the intrinsic $N$-exponent for
resolution-matched work is more plausibly $\approx1.8$--$2.2$.  The honest
bracket is therefore ``$\alpha$ between about $1.3$ measured and $\approx2.2$
implied,'' with the two $N=100$ runs disagreeing on where in that band the truth
lies.

\subsection{What the decomposition says}

The two exponents \eqref{eq:alpha-eq}--\eqref{eq:alpha-rm} bracket the
post-optimization scaling at $\alpha\approx1.3$--$1.8$.  The cost model
\eqref{eq:cost-model} explains the equal-mesh factor of $3.50$ cleanly:
\begin{itemize}
  \item \textbf{Projection-call count} $K\cdot(\mathrm{mBT}+1)$:
    $\dfrac{54{,}536\times2.96}{54{,}154\times2.93}\approx1.02$ --- essentially
    flat.  The iteration count does not grow.
  \item \textbf{Per-call cost} $\mathrm{mPII}\cdot VN$ at fixed $V$:
    $\dfrac{21.15\times100}{13.76\times50}\approx3.08$ --- this is where all the
    growth lives, split into a factor $2$ from $N$ and a factor $1.54$ from the
    rising $\mathrm{mPII}$.
\end{itemize}
Their product ($\approx3.13$) is close to the measured $3.50$; the residual
reflects super-linear per-call overheads (larger $V\times N$ density matrices)
not captured by the leading-order model.  For the equal-mesh comparison the
effective exponent is thus $\approx1.7$--$1.8$, nowhere near~$3$.

The same decomposition applied to the resolution-matched factor of $2.53$ is
what exposes $\alpha_{\text{rm}}=1.34$ as an artifact of the iteration-count
drop rather than favourable per-call scaling:
\begin{itemize}
  \item \textbf{Projection-call count} $K\cdot(\mathrm{mBT}+1)$:
    $\dfrac{16{,}908\times2.48}{54{,}154\times2.93}\approx0.26$ --- it
    \emph{fell} $\sim\!4\times$, driven entirely by the $K$ collapse in the
    finer run.
  \item \textbf{Per-call cost} $\mathrm{mPII}\cdot VN$:
    $\dfrac{32.95\times229{,}216\times100}{13.76\times114{,}144\times50}
    \approx9.6$ --- it \emph{rose} nearly $10\times$ (doubled $V$, doubled $N$,
    $2.4\times$ $\mathrm{mPII}$).
\end{itemize}
Their product ($0.26\times9.6\approx2.5$) recovers the measured factor.  So the
low $1.34$ comes from doing $4\times$ fewer projection calls, not from cheap
per-call scaling --- the per-call cost actually grew steeply.  If that
call-count windfall does not recur at larger $N$ (i.e.\ if $K$ stays flat as it
did at equal mesh), resolution-matched scaling is $\alpha\gtrsim2$.

\begin{remark}[The real driver: per-call projection cost, not iteration count]
Across the three post-optimization runs, $\mathrm{mPII}$ climbs
$13.8\to21\to33$ (Table~\ref{tab:n100-profile}) while $K$ stays flat at equal
mesh.  The growth with $N$ is therefore concentrated in the \emph{cost per
projection call}, consistent with the $\mathcal{O}(\mathrm{mPII}\cdot VN)$
model.  This is exactly the quantity the dual-Newton reformulation of
\S\ref{sec:improvements} attacks.
\end{remark}

\begin{remark}[Uncertainty in $\alpha$]
Even the corrected estimate rests on few points: one $N=50$ run and two $N=100$
runs, all at a single seed.  Run-to-run and host-to-host variance is large (the
same $N=50$ config took $43.9$\,h on a contended host versus $11.6$\,h on a
clean one --- a $3.8\times$ spread from hardware alone).  \textbf{This host
spread ($3.8\times$) exceeds the entire equal-mesh $N$-signal ($3.5\times$)},
so the exponent is only separable from hardware noise if all three runs
executed on comparably clean, uncontended hosts.  The $N=50$ anchor is the
verified clean run; the two $N=100$ runs are \emph{assumed} comparable but this
has not been independently confirmed, and it should be before the bracket is
treated as anything firmer than order-of-magnitude.  Treat
$\alpha\approx1.3$--$1.8$ (measured) / $\approx1.8$--$2.2$ (intrinsic) as an
order-of-magnitude characterisation of the $50\to100$ regime, not a precise law,
and see \S\ref{sec:projections} for why it must not be pushed confidently to
$N=1000$.
\end{remark}

% =========================================================================
\section{Projections to Larger $N$ (Honestly Hedged)}
\label{sec:projections}
% =========================================================================

Anchoring on the measured post-optimization point $T_{50}=11.6$\,h and applying
$T(N)=T_{50}\,(N/50)^{\alpha}$ across the bracket $\alpha\in[1.34,1.81]$ (with
$\alpha=1.5$ as an illustrative midpoint) gives Table~\ref{tab:projections}.
By construction the two $N=100$ columns reproduce the two measured $N=100$ runs
($29.3$\,h at $\alpha=1.34$, $40.6$\,h at $\alpha=1.81$), which is the honest
statement that our two measured points \emph{define} the bracket rather than
test it.

\begin{table}[H]
\centering
\caption{Illustrative Phase~1 wall-time extrapolations on a single clean
machine, anchored on the measured $N=50$ point ($11.6$\,h).  The $N=50$ row and
both $N=100$ entries are \emph{measured}; the $N=200$ and $N=1000$ rows are
extrapolations from only these points and are shown as a range, not a
prediction.  See the caveats below --- especially that $\mathrm{mPII}$ growth
could steepen the effective exponent beyond this bracket.}
\label{tab:projections}
\renewcommand{\arraystretch}{1.3}
\begin{tabular}{rrrr}
\toprule
$N$ & $\alpha=1.34$ & $\alpha=1.5$ (illus.) & $\alpha=1.81$ \\
\midrule
\phantom{0}50   & 11.6\,h \emph{(meas.)} & 11.6\,h \emph{(meas.)} & 11.6\,h \emph{(meas.)} \\
100  & 29.3\,h \emph{(meas.)} & 32.8\,h & 40.6\,h \emph{(meas.)} \\
200  & 74\,h ($\sim$3\,d)   & 93\,h ($\sim$4\,d)   & 143\,h ($\sim$6\,d) \\
1000 & $\sim$640\,h ($\sim$27\,d) & $\sim$1040\,h ($\sim$43\,d) & $\sim$2630\,h ($\sim$110\,d) \\
\bottomrule
\end{tabular}
\end{table}

\noindent Read literally, the measured bracket puts $N=1000$ at roughly
\textbf{one to four months} on a single clean machine; allowing for the steeper
intrinsic (fixed-iteration) exponent $\alpha\approx2$--$2.2$ discussed in
\S\ref{sec:fit} pushes the upper end toward \textbf{$\sim\!6$ months}.  Either
way the honest statement is \textbf{weeks-to-months, not years}.  The practical
conclusions, stated with the appropriate hedging:
\begin{itemize}
  \item $N=50$ is \textbf{routinely feasible} --- measured at half a day on a
    clean host with the seeded initialisation.
  \item $N=100$ is \textbf{measured and feasible} on a single machine: $29.3$\,h
    at matched resolution, $40.6$\,h on the coarser (under-resolved but still
    valid) mesh.  Both produce valid $100$-region partitions.
  \item $N=200$ is a \textbf{multi-day} single-machine run under this bracket
    (roughly $3$--$6$ days); still untested, so treat as an extrapolation.
  \item $N=1000$ is, on this evidence, \textbf{plausibly a weeks-to-months
    single-machine run} rather than an impossibility.  This is an
    \emph{extrapolation from two-to-three points over a $20\times$ range in
    $N$}; it cannot be trusted to better than an order of magnitude.  The
    former ``$\sim\!14$ years / not achievable'' conclusion is
    \textbf{withdrawn} --- it was an artifact of the confounded $N^3$ fit
    (\S\ref{sec:why-old-fit-wrong}), not a property of the algorithm.
\end{itemize}

\begin{remark}[Why $N=1000$ remains genuinely uncertain]
Two-to-three post-optimization points cannot pin a power law over a $20\times$
extrapolation.  The clearest reason the effective exponent could \emph{steepen}
is the very quantity this document identifies as the driver, $\mathrm{mPII}$.
Its growth must be read carefully: of the full $13.8\to21\to33$ climb, only the
first step is $N$-driven ($13.76\to21.15$ at a \emph{fixed} $348\times328$ mesh
as $N$ doubles, a $1.54\times$ rise), while the second ($21.15\to32.95$) is
\emph{mesh}-driven (fixed $N=100$, mesh refined $114$k$\to229$k vertices).  So
the genuinely $N$-attributable growth is the $1.54\times$ per doubling.  If even
that portion is super-constant in $N$, the per-call cost carries more than the
one power of $N$ assumed in \eqref{eq:cost-model}, and $T(N)$ bends upward at
large $N$; a resolution-matched extrapolation compounds it with the
mesh-driven rise as well.  Conversely, the
dual-Newton reformulation (\S\ref{sec:improvements}) would cap $\mathrm{mPII}$
and flatten it.  The point of Table~\ref{tab:projections} is not to assert
$\sim\!43$ days at $N=1000$; it is to show that nothing in the measured data
supports a catastrophic figure, while also refusing to over-commit to an
optimistic one.
\end{remark}

% =========================================================================
\section{Implemented: Serial Optimizations (Changes A/B/C)}
\label{sec:serial-opts}
% =========================================================================

The optimizations below have \emph{landed} (merge \texttt{bae8cd2}); they are
the source of the $\approx\!3.8\times$ constant-factor speedup measured at
$N=50$.  They preserve the optimizer's output --- the $N=50$ partition was
validated bit-identical to the pre-optimization result (contour perimeter
$152.3883843639$, $100\%$ winner-take-all agreement) --- so they reduce the
constant in Eq.~\eqref{eq:cost-model} without changing which factors carry the
$N$-dependence.  The full audit (\#1--\#13, with measured impact tiers) is
recorded in
\texttt{docs/reference/PHASE1\_PGD\_SERIAL\_OPTIMIZATION\_AUDIT.md}.

\paragraph{Change A --- backtracking step warm-start (audit \#1, \#8).}
\texttt{step0} was hard-reset to $1.0$ at the top of every PGD iteration, so
each iteration re-walked the Armijo backtrack from scratch; near convergence
the accepted step decayed to $\sim\!2^{-8}$, i.e.\ $\sim\!8$ rejected trials
per iteration, each a full projection $+$ full energy.  Change~A carries the
last accepted step and starts the next search from
$\min(s_{\max},\,s_{\text{last}}/\rho)$ (Tier~1; Armijo preserved exactly),
and hoists the invariant $\|g\|^2$ out of the trial loop.  Measured effect:
$\mathrm{mBT}$ fell to $1.93$ at $N=50$ (from $6.26$ in the pre-optimization
$N=30$ run) --- the dominant contributor to the speedup, because eliminated
rejected trials remove both their energy \emph{and} their projection
evaluations.

\paragraph{Change B --- projection inner-loop cleanup (audit \#4).}
The constant coupling matrix $C = \|v\|^2(I - J/n)$ and its
$\varepsilon$-regularized solve are now built once per call rather than per
inner iteration; the $\mathcal{O}(VN)$ \texttt{A\_prev = A.copy()} plus
\texttt{np.allclose} stall check is replaced by a scalar consecutive-plateau
test on \texttt{max\_error} (gated by $\texttt{max\_error} < 10\,\mathrm{tol}$
so it never returns an infeasible iterate); the dead
\texttt{convergence\_history} dict is removed.  Result-preserving.  Measured
effect: the projection's share of wall fell from $88$--$90\%$ (pre-opt) to
$80.5\%$ at $N=50$.  (At $N=100$ the share climbs back to $86$--$87\%$, driven
by the higher $\mathrm{mPII}$, not by a regression in Change~B.)

\paragraph{Change C --- post-step gradient reuse (audit \#6).}
The post-step gradient $g_{\text{post}} = \nabla E_\varepsilon(x)$ is
bit-identical to the next iteration's $\nabla E_\varepsilon(x)$ ($x$ is
unmodified between them), so it is carried forward, halving gradient
evaluations.  Bit-identical; $\sim\!1.6\%$ wall.

\begin{remark}[Constant-factor wins do not change the shape of the scaling]
Changes~A/B/C multiply $T$ by a roughly $N$-independent constant
($\approx\!1/3.8$).  They therefore shift every absolute number in
Table~\ref{tab:projections} down but leave the \emph{exponent} to be measured
from the post-optimization runs themselves --- which is exactly what
\S\ref{sec:fit} does, rather than importing a pre-optimization exponent.
\end{remark}

% =========================================================================
\section{What Remains to Scale to $N > 100$}
\label{sec:improvements}
% =========================================================================

Equation~\eqref{eq:cost-model} identifies three independent factors in $T$.
The serial optimizations above shrank the constant; the measured data
(\S\ref{sec:fit}) show the remaining $N$-growth is concentrated in the per-call
projection cost ($\mathrm{mPII}\cdot VN$), \emph{not} in the iteration count.
The durable lever is therefore whatever caps $\mathrm{mPII}$.

\paragraph{1. Further cut re-projection on rejected backtrack steps
(partially done).}
Change~A's Tier~1 warm-start already brought $\mathrm{mBT}$ from the
$4$--$12$ thrashing regime (e.g.\ $6.26$ at pre-opt $N=30$) down to $\sim\!1.9$,
capturing most of the easy gain.  Two stronger options remain, both still
future work:
\begin{itemize}
  \item \textbf{Gradient-step rule without line search} (Barzilai--Borwein or
    spectral step, audit Tier~2): project exactly once per major iteration,
    removing the residual $(\mathrm{mBT}+1)$ multiplier entirely.  Non-monotone;
    needs a Grippo-type safeguard and clamping against the $[10^{-8},1-10^{-8}]$
    box.
  \item \textbf{Lazy projection}: evaluate the unconstrained gradient step
    first and project only when the Armijo condition is about to be accepted.
\end{itemize}
With $\mathrm{mBT}$ already near $2$, the remaining headroom here is modest
($\lesssim\!2\times$) --- the larger structural wins are below.

\paragraph{2. Replace the iterative projection with a dual-Newton closed-form
solve (highest-merit algorithmic fix, audit \#3).}
This is the direct attack on the growing $\mathrm{mPII}$ identified in
\S\ref{sec:fit}.  The coupled sum-to-one $+$ equal-area projection is already
closed-form per inner iteration; only the $A \ge 0$ nonnegativity forces the
$14$--$33$ inner sweeps we now measure.  Dropping the area constraint lets each
vertex row project onto the probability simplex in closed form
($\mathcal{O}(N\log N)$); the $N$ equal-area couplings are then re-introduced
via a dual $y\in\mathbb{R}^N$ solved with a semismooth Newton iteration
($\sim\!5$--$15$ steps).  This replaces the clamp/Sinkhorn sweeps with
$\sim\!10$ Newton steps and gives \emph{exact} feasibility, removing the
$\varepsilon$-floor pathology.  Because $\mathrm{mPII}$ is the quantity that
grows with $N$, capping it is the single change most likely to flatten the
effective exponent; it is the durable fix for $N \ge 100$ and belongs in its
own plan.

\begin{remark}[Inner-solver dual warm-start was ruled out]
An earlier draft listed ``warm-start the projection inner solver'' (carrying
dual multipliers between consecutive PGD projections) as a $\sim\!2\times$
lever.  The serial-optimization audit reclassified this as a \textbf{negative
finding} (audit \#13): the projection's per-call dual state is not a safe warm
start across PGD iterates --- it risks returning an infeasible iterate --- and
the measured profile shows $\mathrm{mPII}$ is only $\sim\!14$ at $N=50$.  The
structural way to attack $\mathrm{mPII}$ (which does grow with $N$) is the
dual-Newton reformulation above, not warm-starting the existing fixed-point
loop.
\end{remark}

\paragraph{3. Parallelise across regions.}
The $\mathcal{O}(N)$ factor in the per-call cost arises from processing each
region's density column independently.  These columns are embarrassingly
parallel: distributing the $V \times N$ density matrix across $P$ cores
reduces per-call cost to $\mathcal{O}(VN/P)$, converting the $N$ factor from
a wall-clock penalty to a throughput one.

\paragraph{4. Iteration count is \emph{not} the bottleneck (corrected).}
An earlier draft claimed $K\propto N^2$ and listed reducing that exponent as a
priority.  The measured post-optimization data refute the premise: at equal
mesh $K$ is flat ($54{,}154\to54{,}536$ from $N=50$ to $N=100$;
Table~\ref{tab:n100-profile}).  There is no $N^2$ iteration-count growth to
remove.  Better initialisation (the seeded Voronoi approach) and adaptive
$\varepsilon$ remain worthwhile for robustness, but they are not the lever for
$N$-scaling --- the per-call projection cost is.  Confirming the flat-$K$ trend
at a third $N$ (e.g.\ $N=200$) with the seeded init would further solidify the
corrected picture.

\paragraph{5. FEM assembly is not the bottleneck.}
Assembly is $<1\%$ of wall time at current mesh sizes; it becomes relevant only
above $\sim\!10^5$--$10^6$ vertices.

\begin{remark}[Cluster scaling]
Running multiple seeds in parallel on a cluster (e.g.\ Pelle) does not reduce
the per-seed wall time but gives the ensemble result sooner.  For $N=100$
($\sim\!1.2$--$1.7$ days per seed, measured) this compresses the effective
turnaround when several seeds are needed.  For $N=1000$ --- a weeks-to-months
per-seed run under the current bracket --- an ensemble shortens turnaround but
does not change the per-seed cost; the dual-Newton projection is the lever that
does.
\end{remark}

% =========================================================================
\section*{Summary table}
% =========================================================================

\begin{center}
\renewcommand{\arraystretch}{1.35}
\begin{tabular}{p{2.8cm}p{4.6cm}p{5.1cm}}
\toprule
\textbf{Factor} & \textbf{Current scaling (post-opt, measured)} & \textbf{Path to improvement} \\
\midrule
Cost per projection call
  & $\mathcal{O}(\mathrm{mPII}\cdot VN)$; $\mathrm{mPII}$ grows
    $13.8\to21\to33$ across the measured runs; projection $80$--$87\%$ of wall
  & Dual-Newton closed-form simplex projection (audit \#3) --- the lever;
    parallelise regions ($P\times$) \\
\# projection calls
  & $K\cdot(\mathrm{mBT}+1)$; $\mathrm{mBT}=1.93$ after Change~A (was $4$--$12$)
  & BB/spectral step or lazy projection: remove residual $\mathrm{mBT}+1$
    multiplier \\
Iteration count $K$
  & \textbf{Flat}, not $\sim\!N^2$: $54{,}154\to54{,}536$ ($N{=}50\to100$,
    equal mesh)
  & Not a scaling bottleneck; better init aids robustness only \\
Overall exponent $\alpha$
  & Measured $\approx1.3$--$1.8$ ($N{=}50\to100$); intrinsic (fixed-$K$)
    $\approx1.8$--$2.2$ --- the $1.3$ end is inflated by a one-off $4\times$
    iteration-count drop.  Either way \emph{not} $3.1$
  & Cap $\mathrm{mPII}$ via the dual-Newton projection to flatten it \\
Serial opts (A/B/C)
  & Landed (\texttt{bae8cd2}); $\approx\!3.8\times$, output-identical
  & Done --- constant-factor only \\
$N=50$ (measured)
  & $11.6$\,h clean host, valid partition
  & Routinely feasible on one machine \\
$N=100$ (measured, both valid)
  & $29.3$\,h (matched mesh) / $40.6$\,h (coarser mesh)
  & Feasible now; cluster for multi-seed \\
$N=1000$ (extrapolated)
  & $\sim\!43$\,d ($\alpha{=}1.5$) to $\sim\!190$\,d ($\alpha{=}2$);
    \emph{weeks-to-months}, uncertain to $>\!1$ order of magnitude
  & Dual-Newton projection $+$ parallelism; ``14\,yr / not achievable''
    retracted \\
\bottomrule
\end{tabular}
\end{center}

\bibliographystyle{plain}
\bibliography{../shared/references}

\end{document}
